\documentclass[11pt,a4paper]{article}
\usepackage{konspekt}
\renewcommand{\kheadleft}{Теория вероятностей и статистика}
\hypersetup{pdftitle={Конспект: теория вероятностей и статистика},pdfauthor={Конспект теории}}

\begin{document}

\makecover{Теория вероятностей\\и статистика}
{\begin{tikzpicture}[scale=1.15]
  \foreach \x/\h in {-3/0.06,-2.5/0.14,-2/0.32,-1.5/0.62,-1/1.02,-0.5/1.38,0/1.52,0.5/1.36,1/1.0,1.5/0.64,2/0.3,2.5/0.15,3/0.05}
    \fill[korange!22] (\x-0.24,0) rectangle (\x+0.24,{1.6*\h});
  \fill[kviolet!14,domain=-1:1,samples=60] (-1,0) -- plot (\x,{2.45*exp(-\x*\x/2)}) -- (1,0) -- cycle;
  \draw[kviolet,line width=1.6pt,domain=-3.6:3.6,samples=120] plot (\x,{2.45*exp(-\x*\x/2)});
  \draw[axis] (-4,0) -- (4.2,0);
  \foreach \x/\l in {-2/$-2\sigma$,-1/$-\sigma$,0/$\E X$,1/$+\sigma$,2/$+2\sigma$}
    \draw (\x,0.07) -- (\x,-0.07) node[below,font=\footnotesize] {\l};
  \draw[aux] (0,0) -- (0,2.45);
  \begin{scope}[shift={(4.6,2.6)},rotate=12]
    \draw[line width=1pt,rounded corners=4pt,fill=white] (0,0) rectangle (1.3,1.3);
    \foreach \p in {(0.33,0.33),(0.97,0.97),(0.65,0.65)} \fill[korange] \p circle (0.11);
  \end{scope}
  \begin{scope}[shift={(5.6,1.2)},rotate=-14]
    \draw[line width=1pt,rounded corners=4pt,fill=white] (0,0) rectangle (1.3,1.3);
    \foreach \p in {(0.33,0.33),(0.97,0.97),(0.33,0.97),(0.97,0.33),(0.33,0.65),(0.97,0.65)} \fill[kviolet] \p circle (0.11);
  \end{scope}
  \begin{scope}[shift={(-6.6,0.4)},scale=0.55]
    \draw[axis] (0,0) -- (4.6,0); \draw[axis] (0,0) -- (0,4.2);
    \foreach \x/\y in {0.5/0.9,0.9/0.7,1.2/1.5,1.6/1.3,2/2.1,2.3/1.7,2.7/2.6,3/2.3,3.4/3.1,3.8/3.0,1.9/1.0,3.1/3.4}
      \fill[kviolet] (\x,\y) circle (0.09);
    \draw[korange,line width=1.2pt] (0.2,0.45) -- (4.3,3.6);
  \end{scope}
\end{tikzpicture}}

\maketoc

\section{Ковариация и корреляция}

\subsection{Свойства математического ожидания и дисперсии}
Для любых случайных величин $X$, $Y$ и чисел $a$, $b$:
\[
\kbox{\E(aX+bY)=a\E(X)+b\E(Y),\qquad \D(aX+bY)=a^2\D(X)+b^2\D(Y)+2ab\cov(X,Y).}
\]

\begin{zam}
Если $X$ и $Y$ независимы, то $\cov(X,Y)=0$ и $\D(aX+bY)=a^2\D(X)+b^2\D(Y)$.
В частности, $\D(X+Y)=\D(X-Y)=\D(X)+\D(Y)$: \akc{при вычитании дисперсии складываются}.
\end{zam}

\subsection{Ковариация и коэффициент корреляции}

\begin{opr}[ковариация]
\term{Ковариацией} случайных величин $X$ и $Y$ называется число
\[
\cov(X,Y)=\E\bigl[(X-\E X)(Y-\E Y)\bigr]=\E(XY)-\E(X)\E(Y).
\]
\end{opr}

\begin{opr}[коэффициент корреляции]
\[
\rho(X,Y)=\frac{\cov(X,Y)}{\sigma(X)\,\sigma(Y)}=\frac{\cov(X,Y)}{\sqrt{\D(X)\D(Y)}},\qquad -1\le\rho(X,Y)\le1.
\]
\end{opr}

Отсюда $\cov(X,Y)=\rho(X,Y)\sqrt{\D(X)\D(Y)}$ и $\E(XY)=\cov(X,Y)+\E(X)\E(Y)$.

\begin{kob}{Связь с независимостью}
Если $X$ и $Y$ независимы, то $\cov(X,Y)=0$ и $\rho(X,Y)=0$.
Следовательно, если $\rho(X,Y)\ne0$, то $X$ и $Y$ \textbf{зависимы}.
\end{kob}

\begin{kvb}{Индикаторы событий}
Пусть $A$, $B$ — индикаторы событий (принимают значение $1$, если событие произошло, и $0$ иначе).
Тогда $\E(A)=\P(A)$, $\D(A)=\P(A)\bigl(1-\P(A)\bigr)$, а произведение $AB=1$ только если $A=1$ и $B=1$, поэтому
\[
\E(AB)=\P(AB)=\cov(A,B)+\P(A)\P(B).
\]
\end{kvb}

\section{Реальные (выборочные) модели}

\begin{opr}[выборка]
\term{Выборкой}, связанной с совместным наблюдением случайных величин $X$ и $Y$, будем называть $N$ пар значений
\[
(x_1,y_1),\ (x_2,y_2),\ \dots,\ (x_N,y_N)
\]
таких, что $x_1,\dots,x_N$ имеют такое же распределение, что и $X$, а $y_1,\dots,y_N$ — такое же распределение, что и $Y$.
\end{opr}

Каждой характеристике вероятностной модели соответствует её выборочный аналог:

\begin{table}[H]
\centering
\renewcommand{\arraystretch}{2.1}
\begin{tabular}{>{\centering\arraybackslash}m{4.2cm} >{\centering\arraybackslash}m{9.6cm}}
\toprule
{\hfont\color{kviolet}Вероятностная модель} & {\hfont\color{kviolet}Реальная (выборочная) модель}\\
\midrule
$\E(X)$ & $\bar x=\dfrac{x_1+\dots+x_N}{N}$\\
$\D(X)$ & $S_x^2=\dfrac{(x_1-\bar x)^2+\dots+(x_N-\bar x)^2}{N}=\overline{x^2}-\bar x^{\,2}$\\
$\sigma(X)$ & $S_x=\sqrt{S_x^2}$\\
$\cov(X,Y)$ & $\begin{aligned}C_{xy}&=\dfrac{(x_1-\bar x)(y_1-\bar y)+\dots+(x_N-\bar x)(y_N-\bar y)}{N}\\[-2pt]&=\overline{xy}-\bar x\,\bar y\end{aligned}$\\[14pt]
$\rho(X,Y)$ & $R_{xy}=\dfrac{C_{xy}}{S_xS_y}$\\
\bottomrule
\end{tabular}
\caption{Характеристики вероятностной модели и их выборочные аналоги}
\end{table}

\subsection{Диаграмма рассеивания}
Зависимость между величинами удобно изучать с помощью \term{диаграммы рассеивания} (облака рассеивания):
каждая пара $(x_i,y_i)$ изображается точкой на координатной плоскости. Если точки группируются вдоль прямой,
зависимость близка к линейной, и её можно описать уравнением $y=ax+b$.

\begin{figure}[H]
\centering
\begin{tikzpicture}[scale=1]
  \draw[axis] (0,0) -- (8,0) node[below] {$x$};
  \draw[axis] (0,0) -- (0,4.8) node[left] {$y$};
  \foreach \x/\y in {0.8/0.6,1.2/1.2,1.6/0.9,2.1/1.5,2.5/1.3,2.9/2.0,3.3/1.6,3.7/2.4,4.1/2.1,4.6/2.9,6.0/3.3,6.3/3.7,6.8/3.4,4.0/2.9,2.0/0.9}
     \fill[kviolet] (\x,\y) circle (2pt);
  \draw[korange,line width=1.2pt] (0.4,0.55) -- (7.5,4.05) node[above left] {$y=ax+b$};
  \fill[kviolet] (5.2,2.0) circle (2pt);
  \draw[korange,dashed,line width=0.9pt] (5.2,2.0) -- (5.2,2.92);
  \fill[korange] (5.2,2.92) circle (1.6pt);
  \draw[korange,line width=0.4pt] (5.25,2.45) -- (5.9,1.6) node[below right,font=\footnotesize,inner sep=1pt] {$ax_i+b-y_i$};
  \draw[thin aux] (5.2,0) -- (5.2,2.0);
  \node[below,font=\footnotesize] at (5.2,0) {$x_i$};
\end{tikzpicture}
\caption{Диаграмма рассеивания и прямая регрессии}
\end{figure}

\subsection{Метод наименьших квадратов для линейной регрессии}
Подбираем прямую $y=ax+b$ так, чтобы сумма квадратов отклонений была минимальной:
\[
S(a,b)=(ax_1+b-y_1)^2+(ax_2+b-y_2)^2+\dots+(ax_N+b-y_N)^2\longrightarrow\min.
\]
Найти минимум можно двумя способами: через частные производные или \emph{выделением полных квадратов}. Разберём второй.

Раскроем скобки:
\begin{multline*}
S(a,b)=(y_1^2+\dots+y_N^2)+a^2(x_1^2+\dots+x_N^2)+Nb^2\\
-2a(x_1y_1+\dots+x_Ny_N)-2b(y_1+\dots+y_N)+2ab(x_1+\dots+x_N).
\end{multline*}
Разделим на $N$:
\[
\frac{S(a,b)}{N}=\overline{y^2}+a^2\,\overline{x^2}+b^2-2a\,\overline{xy}-2b\,\bar y+2ab\,\bar x.
\]
Выделим полный квадрат по $b$, используя равенства
\[
\overline{x^2}=S_x^2+\bar x^{\,2},\qquad \overline{y^2}=S_y^2+\bar y^{\,2},\qquad \overline{xy}=C_{xy}+\bar x\,\bar y:
\]
\[
\frac{S(a,b)}{N}=(a\bar x+b-\bar y)^2+a^2S_x^2-2aC_{xy}+S_y^2.
\]
Теперь выделим полный квадрат по $a$:
\[
\frac{S(a,b)}{N}=(a\bar x+b-\bar y)^2+\Bigl(aS_x-\frac{C_{xy}}{S_x}\Bigr)^2\underbrace{-\frac{C_{xy}^2}{S_x^2}+S_y^2}_{\const}.
\]
Минимум достигается, когда оба квадрата равны нулю:
\[
\begin{cases}
a\bar x+b-\bar y=0,\\[2pt]
aS_x-\dfrac{C_{xy}}{S_x}=0.
\end{cases}
\]

\begin{kvb}{Коэффициенты прямой регрессии}
\[
a=\frac{C_{xy}}{S_x^2}=R_{xy}\frac{S_y}{S_x},\qquad b=\bar y-a\bar x=\bar y-R_{xy}\frac{S_y}{S_x}\,\bar x.
\]
\end{kvb}

\begin{zam}
Прямая регрессии всегда проходит через точку $(\bar x,\bar y)$ — «центр» облака рассеивания.
\end{zam}

\section{Закон больших чисел}

\subsection{Неравенство Маркова}

\begin{kob}{Неравенство Маркова}
Пусть случайная величина $X\ge0$. Тогда для любого $\varepsilon>0$
\[
\P(X\ge\varepsilon)\le\frac{\E(X)}{\varepsilon}.
\]
\end{kob}

\begin{proof}
Пусть $X$ принимает значения $x_1,\dots,x_n$ с вероятностями $p_1,\dots,p_n$, причём $x_1,\dots,x_k<\varepsilon$,
а $x_{k+1},\dots,x_n\ge\varepsilon$. Тогда
\[
\E(X)=x_1p_1+\dots+x_np_n=S_1+S_2,\qquad S_1=x_1p_1+\dots+x_kp_k,\quad S_2=x_{k+1}p_{k+1}+\dots+x_np_n.
\]
Так как $X\ge0$, то $S_1\ge0$, $S_2\ge0$, и
\[
\E(X)\ge S_2\ge\varepsilon(p_{k+1}+\dots+p_n)=\varepsilon\,\P(X\ge\varepsilon).
\]
Отсюда $\P(X\ge\varepsilon)\le\dfrac{\E(X)}{\varepsilon}$.
\end{proof}

\subsection{Неравенство Чебышёва}

\begin{kob}{Неравенство Чебышёва}
Пусть случайная величина $X$ имеет математическое ожидание $\E(X)$ и дисперсию $\D(X)$. Тогда для любого $\varepsilon>0$
\[
\P\bigl(|X-\E(X)|\ge\varepsilon\bigr)\le\frac{\D(X)}{\varepsilon^2}.
\]
\end{kob}

\begin{proof}
Рассмотрим $Y=(X-\E(X))^2\ge0$; заметим, что $\E(Y)=\D(X)$. По неравенству Маркова
\[
\P(Y\ge\varepsilon^2)\le\frac{\E(Y)}{\varepsilon^2}
\;\Longrightarrow\;
\P\bigl((X-\E(X))^2\ge\varepsilon^2\bigr)\le\frac{\D(X)}{\varepsilon^2}
\;\Longrightarrow\;
\P\bigl(|X-\E(X)|\ge\varepsilon\bigr)\le\frac{\D(X)}{\varepsilon^2}.\qedhere
\]
\end{proof}

Эквивалентная форма (для противоположного события):
\[
\kbox{\P\bigl(|X-\E(X)|<\varepsilon\bigr)\ge1-\frac{\D(X)}{\varepsilon^2}.}
\]

\begin{zam}
Неравенство Чебышёва даёт, как правило, более точную оценку, чем неравенство Маркова,
так как учитывает разброс величины (дисперсию).
\end{zam}

\subsection{Следствия}
\begin{enumerate}
\item Пусть $X\ge0$ и $\E(X)=a$. Тогда для $k\in\NN$
\[
\P(X\ge ka)\le\frac{a}{ka}=\frac1k.
\]
\item \textbf{Правило $k$ сигм.} Подставим $\varepsilon=k\sigma$ в неравенство Чебышёва:
\[
\P(|X-\E X|\ge2\sigma)\le\frac{\sigma^2}{4\sigma^2}=\frac14,\qquad
\P(|X-\E X|\ge3\sigma)\le\frac19,\qquad
\P(|X-\E X|\ge4\sigma)\le\frac1{16}.
\]
\end{enumerate}

\begin{kob}{Правило трёх сигм}
Для нормально распределённой величины оценка намного точнее: с вероятностью $\approx0{,}997$
\[
X\in\bigl[\E(X)-3\sigma;\ \E(X)+3\sigma\bigr].
\]
\end{kob}

\begin{figure}[H]
\centering
\begin{tikzpicture}[xscale=1.25,yscale=1.25]
  \def\f(#1){2.4*exp(-(#1)*(#1)/2)}
  \fill[kviolet!8,domain=-3:3,samples=80] (-3,0) -- plot (\x,{\f(\x)}) -- (3,0) -- cycle;
  \fill[kviolet!18,domain=-2:2,samples=80] (-2,0) -- plot (\x,{\f(\x)}) -- (2,0) -- cycle;
  \fill[kviolet!30,domain=-1:1,samples=60] (-1,0) -- plot (\x,{\f(\x)}) -- (1,0) -- cycle;
  \draw[kviolet,line width=1.3pt,domain=-3.7:3.7,samples=120] plot (\x,{\f(\x)});
  \draw[axis] (-4,0) -- (4.1,0);
  \foreach \x/\l in {-3/$\E X-3\sigma$,-2/$-2\sigma$,-1/$-\sigma$,0/$\E X$,1/$+\sigma$,2/$+2\sigma$,3/$\E X+3\sigma$}
    \draw (\x,0.06) -- (\x,-0.06) node[below,font=\footnotesize] {\l};
  \node[font=\small] at (0,0.8) {$68\%$};
  \draw[korange,<->,line width=0.8pt] (-2,3.0) -- (2,3.0) node[midway,above,font=\footnotesize,korange] {$\approx95\%$};
  \draw[korange,<->,line width=0.8pt] (-3,3.6) -- (3,3.6) node[midway,above,font=\footnotesize,korange] {$\approx99{,}7\%$};
  \foreach \x in {-3,-2,2,3} \draw[thin aux] (\x,0) -- (\x,{\x*\x<5?3.0:3.6});
\end{tikzpicture}
\caption{Правило трёх сигм для нормального распределения}
\end{figure}

\section{Схема Бернулли (биномиальное распределение)}

Проводится $n$ независимых одинаковых испытаний, в каждом из которых возможен «успех» с вероятностью $p$
или «неудача» с вероятностью $q=1-p$.

\begin{opr}[формула Бернулли]
Вероятность того, что в $n$ испытаниях произойдёт ровно $k$ успехов:
\[
P_n(S=k)=C_n^k\,p^k q^{\,n-k},
\]
где $n$ — число испытаний, $k$ — число успехов, $p$ — вероятность успеха, $q=1-p$ — вероятность неудачи.
\end{opr}

Пусть $S$ — случайная величина, обозначающая число успехов в испытаниях Бернулли. Её закон распределения называется
\term{биномиальным}:
\[
S\sim\begin{pmatrix}
0 & 1 & 2 & \dots & k & \dots & n\\
q^n & C_n^1pq^{n-1} & C_n^2p^2q^{n-2} & \dots & C_n^kp^kq^{n-k} & \dots & p^n
\end{pmatrix}.
\]

\begin{figure}[H]
\centering
\begin{tikzpicture}
\begin{axis}[kplot,width=11cm,height=5.4cm,ymin=0,ymax=0.3,xmin=-0.7,xmax=10.7,
  xtick={0,...,10},ytick={0,0.1,0.2},yticklabel style={/pgf/number format/.cd,use comma,fixed},
  xlabel={$k$},ylabel={$P_n(S=k)$},ymajorgrids,
  title={\small Биномиальное распределение, $n=10$, $p=0{,}3$}]
\addplot[ybar,bar width=14pt,fill=kviolet!70,draw=kviolet] coordinates
 {(0,0.0282)(1,0.1211)(2,0.2335)(3,0.2668)(4,0.2001)(5,0.1029)(6,0.0368)(7,0.0090)(8,0.0014)(9,0.0001)(10,0.0000)};
\end{axis}
\end{tikzpicture}
\caption{Биномиальное распределение}
\end{figure}

\begin{teo}
Пусть $S$ — случайная величина, имеющая биномиальное распределение. Тогда
\[
\E(S)=np,\qquad \D(S)=npq,
\]
где $n$ — число испытаний, $p$ — вероятность успеха в отдельном испытании, $q$ — вероятность неудачи.
\end{teo}

\begin{proof}
Пусть $S_1$ — число успехов в 1-м испытании:
\[
S_1\sim\begin{pmatrix}0&1\\q&p\end{pmatrix}\;\Longrightarrow\;\E(S_1)=p,\qquad
\D(S_1)=\E(S_1^2)-\E^2(S_1)=p-p^2=p(1-p)=pq.
\]
Аналогично для $S_2,S_3,\dots,S_n$. Так как $S=S_1+S_2+\dots+S_n$, то
\[
\E(S)=\E(S_1)+\dots+\E(S_n)=np,\qquad \D(S)=\D(S_1)+\dots+\D(S_n)=npq,
\]
где дисперсии складываются, \emph{так как испытания независимы}.
\end{proof}

\subsection{Частота появления события}
Рассмотрим серию из $N$ одинаковых опытов, в каждом из которых может произойти событие $A$ с вероятностью $p$.
Пусть $S$ — число появлений события $A$; тогда $\E(S)=Np$, $\D(S)=Npq$.

Рассмотрим случайную величину $Y=\dfrac SN$ — \term{частоту} появления успеха:
\[
\E(Y)=\E\Bigl(\frac SN\Bigr)=\frac1N\,pN=p,\qquad
\D(Y)=\D\Bigl(\frac SN\Bigr)=\frac1{N^2}\,Npq=\frac{pq}{N}\xrightarrow[N\to\infty]{}0.
\]
По неравенству Чебышёва для любого $\varepsilon>0$:
\[
\P(|Y-p|\ge\varepsilon)\le\frac{\D(Y)}{\varepsilon^2}=\frac{pq}{N\varepsilon^2}\xrightarrow[N\to\infty]{}0,\qquad
\P(|Y-p|<\varepsilon)\ge1-\frac{pq}{N\varepsilon^2}\xrightarrow[N\to\infty]{}1.
\]

\begin{kob}{Теорема Бернулли}
С вероятностью, сколь угодно близкой к $1$, можно утверждать, что при достаточно большом числе испытаний $N$
частота случайного события будет сколь угодно мало отличаться от его вероятности.
\end{kob}

\begin{fbx}{Полезные формулы для числа успехов $X$}
\[
\P(|X-np|\ge\varepsilon)\le\frac{npq}{\varepsilon^2},\qquad
\P\Bigl(\Bigl|\frac XN-p\Bigr|\ge\varepsilon\Bigr)\le\frac{pq}{N\varepsilon^2}.
\]
Так как $pq=p(1-p)\le\frac14$, вторую оценку можно заменить на $\dfrac{1}{4N\varepsilon^2}$, если $p$ неизвестно.
\end{fbx}

\section{Теорема Чебышёва}

\begin{kob}{Теорема Чебышёва}
С вероятностью, сколь угодно близкой к $1$, можно утверждать, что при достаточно большом числе независимых опытов $N$
среднее арифметическое полученных значений случайной величины будет сколь угодно мало отличаться от её математического ожидания.
\end{kob}

\begin{proof}
Пусть $\E(X)=m$, $\D(X)=d$, а $X_1,\dots,X_N$ — значения $X$ в независимых опытах.
Рассмотрим среднее $\bar X=\dfrac{X_1+X_2+\dots+X_N}{N}$. Тогда
\[
\E(\bar X)=\frac1N\bigl(\E(X_1)+\dots+\E(X_N)\bigr)=\frac{Nm}{N}=m,
\]
\[
\D(\bar X)=\frac1{N^2}\bigl(\D(X_1)+\dots+\D(X_N)\bigr)=\frac{Nd}{N^2}=\frac dN\xrightarrow[N\to\infty]{}0
\]
(дисперсии складываются в силу независимости). По неравенству Чебышёва для любого $\varepsilon>0$:
\[
\P(|\bar X-m|\ge\varepsilon)\le\frac{d}{N\varepsilon^2}\;\Longrightarrow\;
\P(|\bar X-m|<\varepsilon)\ge1-\frac{d}{N\varepsilon^2}\xrightarrow[N\to\infty]{}1.\qedhere
\]
\end{proof}

\begin{zam}
Теорема Бернулли — частный случай теоремы Чебышёва: частота $\frac SN$ есть среднее арифметическое индикаторов
успеха $S_1,\dots,S_N$, у которых $m=p$, $d=pq$.
\end{zam}

\section{Геометрическое распределение}

Испытания Бернулли проводятся \emph{до первого успеха}. Пусть $X$ — номер испытания, в котором впервые произошёл успех.

\begin{opr}[геометрическое распределение]
\[
\P(X=k)=p\,q^{k-1},\qquad k=1,2,3,\dots,
\]
где $p$ — вероятность успеха, $q=1-p$, испытания независимы.
\[
X\sim\begin{pmatrix}1&2&3&\dots&k&\dots\\ p&pq&pq^2&\dots&pq^{k-1}&\dots\end{pmatrix}
\]
\end{opr}

Проверим, что сумма вероятностей равна $1$ (геометрическая прогрессия):
\[
p+pq+pq^2+\dots+pq^{k-1}+\dots=p(1+q+q^2+\dots)=p\cdot\frac1{1-q}=1.
\]

\begin{figure}[H]
\centering
\begin{tikzpicture}
\begin{axis}[kplot,width=11cm,height=5.4cm,ymin=0,ymax=0.45,xmin=0.3,xmax=10.7,
  xtick={1,...,10},ytick={0,0.1,0.2,0.3,0.4},yticklabel style={/pgf/number format/.cd,use comma,fixed},
  xlabel={$k$},ylabel={$\P(X=k)$},ymajorgrids,
  title={\small Геометрическое распределение, $p=0{,}4$}]
\addplot[ybar,bar width=14pt,fill=korange!80,draw=korange] coordinates
 {(1,0.4)(2,0.24)(3,0.144)(4,0.0864)(5,0.0518)(6,0.0311)(7,0.0187)(8,0.0112)(9,0.0067)(10,0.0040)};
\end{axis}
\end{tikzpicture}
\caption{Геометрическое распределение}
\end{figure}

\subsection{Математическое ожидание}
\[
\E(X)=1\cdot p+2\cdot pq+3\cdot pq^2+\dots+k\,pq^{k-1}+\dots
\]
Разобьём каждое слагаемое: $2pq=pq+pq$, $3pq^2=pq^2+2pq^2$, $4pq^3=pq^3+3pq^3$ и т.\,д. Тогда
\[
\E(X)=\underbrace{(p+pq+pq^2+\dots)}_{=1}+q\underbrace{(p+2pq+3pq^2+\dots)}_{=\E(X)}=1+q\E(X).
\]
Отсюда $\E(X)(1-q)=1$, то есть
\[
\kbox{\E(X)=\frac1p.}
\]

\subsection{Дисперсия}
$\D(X)=\E(X^2)-\E^2(X)$. Найдём
\[
\E(X^2)=1\cdot p+4pq+9pq^2+16pq^3+\dots
\]
Разобьём: $4pq=pq+3pq$, $9pq^2=4pq^2+5pq^2$, $16pq^3=9pq^3+7pq^3$, \dots{} (здесь $3=2\cdot2-1$, $5=2\cdot3-1$, $7=2\cdot4-1$,
т.\,е. $k^2=(k-1)^2+(2k-1)$). Тогда
\begin{align*}
\E(X^2)&=q\underbrace{(p+4pq+9pq^2+\dots)}_{\E(X^2)}+\underbrace{(p+3pq+5pq^2+7pq^3+\dots)}_{\E(2X-1)}\\
&=q\E(X^2)+2\E(X)-1=q\E(X^2)+\frac2p-1.
\end{align*}
Значит, $p\E(X^2)=\dfrac2p-1$ и $\E(X^2)=\dfrac2{p^2}-\dfrac1p$. Окончательно
\[
\D(X)=\frac2{p^2}-\frac1p-\frac1{p^2}=\frac{1-p}{p^2},\qquad \kbox{\D(X)=\frac{q}{p^2}.}
\]

\begin{itog}
Для геометрического распределения:
\[
\P(X=k)=pq^{k-1},\qquad \E(X)=\frac1p,\qquad \D(X)=\frac q{p^2},\qquad \P(X\le k)=1-q^k.
\]
\end{itog}

\end{document}
